\section{From stitching to federation}
\label{sec:ch01-stitching-to-federation}

When Facebook open-sourced GraphQL in 2015, the specification described a
language for asking one server for data~\autocite{byron2015graphql}. It had
nothing to say about how several teams should build that server together,
because at the time that was not the problem anybody had. The problem arrived
about two years later, and the industry has been working on it ever since. The
history is worth a few pages, because both of the answers people tried are
still in use, and the difference between them is the difference between a
gateway you maintain and a gateway you configure.

\subsection*{Stitching}
\index{schema stitching}

Apollo published the first answer in September 2017. Schema stitching was, in
Sashko Stubailo's description, \enquote{the idea that you can take two or more
GraphQL schemas, and merge them into one endpoint that can pull data from all
of them}~\autocite{stubailo2017stitching}. It shipped as a stable feature that
October in \code{graphql-tools} 2.0, with three moving
parts~\autocite{novikov2017tools2}: introspect a remote API and turn it into a
local schema object, merge several such schemas into one, and write resolvers
by hand to link types that live in different backends.

That last part is the one to hold on to. If \code{Order} lived in one service
and \code{User} in another, the code that taught the graph how to get from an
order to its customer was written in the gateway, by whoever owned the gateway.

Stitching had a real virtue, and it is the reason its descendants are still
around. It asked nothing of the services being stitched. Any spec-compliant
GraphQL API could be absorbed without knowing it had been absorbed, which is
exactly what you want when you are aggregating things you do not control, such
as a vendor's API or a service belonging to a team that has no interest in your
architecture.

What it did not fix was the queue from section~\ref{sec:ch01-org-chart}. It
moved it. Instead of six teams editing one schema file, six teams now needed
changes in one gateway repository, and the gateway's resolvers encoded
assumptions about types those teams owned. Apollo's 2018 release added a
transform layer, which made the glue more capable and also gave it more surface
area to be wrong about~\autocite{stubailo2018nextgen}.

The sharpest contemporary statement of the problem came from Marc-Andr\'e
Giroux in October 2018, six months before federation existed, which is part of
why I trust it more than the migration stories that came later. He observed
that a large stitched schema ends up with \enquote{a lot of `glue' code at the
gateway itself, defining all the important relations there}, that
responsibility drifts until the gateway becomes \enquote{responsible of
everything}, and then asked the question that has never really been answered on
stitching's own terms: \enquote{what is the point of having decentralized
schemas when in reality, we need a centralized point to merge and extend the
schemas?} He also noticed the cost that shows up in a developer's week rather
than an architecture review, which is that linking a new field to another
service now takes two steps instead of one, once at the service and again at
the gateway~\autocite{giroux2018stitching}.

Expedia's account of migrating off stitching names the operational version of
the same thing: their gateway code kept getting more complex, and there was
\enquote{no way to determine the `true schema' without running the
gateway}~\autocite{isquick2021expedia}. A composed schema that only exists at
runtime cannot be reviewed, diffed, or checked in CI, which means nothing can
tell you that today's deploy just broke a client.

\subsection*{Federation}
\index{Apollo Federation!origins}

Apollo announced Federation on 1 May 2019, and was direct about what it was
for: replacing stitching and solving \enquote{pain points such as coordination,
separation of concerns, and brittle gateway code}~\autocite{baxley2019federation}.
The design move was to invert where the linking logic lives. Rather than a
gateway that knows how to join services, each service declares which types it
contributes to and how it can be looked up, and a build step works out the
joins. In the announcement's words, \enquote{you compose a graph declaratively
from within your schema instead of writing imperative schema stitching code}.

The same post contains the sentence that this book is arguably a long footnote
to: \enquote{Often no single team controls every aspect of an important type
like a User or Product, so the definition of these types should be distributed
across teams and codebases, rather than centralized.} That is the Conway
argument from the previous section, made by the vendor, three years before most
of us were forced to agree with it.

\index{entity}\index{key directive}%
The mechanism is an \emph{entity}: a type that more than one service can
contribute fields to, with a declared key that identifies one of them. Two
subgraphs describing the same product look roughly like this:

\begin{graphqlsdl}
# Sketch. Catalog owns the product's identity and description.
type Product @key(fields: "id") {
  id: ID!
  title: String!
}

# In a different service, a different repository, a different deploy.
type Product @key(fields: "id") {
  id: ID!
  price: Money!
}
\end{graphqlsdl}

Neither service knows about the other. Both claim \code{Product}, and both agree
on how a product is identified. Everything federation does follows from that:
a build step checks the claims are compatible and produces one schema, and at
request time something has to fetch the pieces and stitch the result together.
Chapter~\ref{ch:the-federation-model} covers the directives properly and
chapter~\ref{ch:how-a-federated-query-actually-runs} follows the HTTP traffic,
so I will leave both alone here. The shape is what matters for now.

\begin{figure}[htbp]
  \centering
  % The shape of a federated graph: independent services, one composed schema,
% one router in front. Referenced from chapter 01. Bare tikzpicture; the figure
% environment, caption and label live at the call site.
\begin{tikzpicture}[
    font=\small,
    box/.style={draw, rounded corners=2pt, minimum width=24mm,
                minimum height=8mm, align=center},
    router/.style={draw, thick, minimum width=40mm, minimum height=10mm,
                   align=center},
    super/.style={draw, dashed, minimum width=38mm, minimum height=11mm,
                  align=center},
    consumer/.style={draw, dotted, minimum width=30mm, minimum height=8mm,
                     align=center},
    flow/.style={-{Stealth[length=2mm]}, thick},
    build/.style={-{Stealth[length=2mm]}, dashed},
    note/.style={font=\scriptsize\itshape, align=center}
  ]

  \node[consumer] (client) at (0,3.1) {one client query};
  \node[router]   (router) at (0,1.7) {router};
  \node[super]    (super)  at (4.9,1.7) {composed\\supergraph schema};

  \node[box] (sg1) at (-3.1,-0.4) {Catalog};
  \node[box] (sg2) at ( 0,  -0.4) {Orders};
  \node[box] (sg3) at ( 3.1,-0.4) {Pricing};

  \draw[flow]  (client) -- (router);
  \draw[build] (super)  -- (router);

  \foreach \s in {sg1,sg2,sg3}{%
    \draw[flow] (router.south) -- (\s.north);
  }

  \node[note, anchor=north] at (4.9,1.05) {built at deploy time};
  \node[note, anchor=north] at (0,-1.15)
    {three services, three teams, three deploy schedules};

\end{tikzpicture}

  \caption{Independent services, one schema composed from them at deploy time,
  one router serving clients. Compare this with
  figure~\ref{fig:ch01-one-schema-many-teams}: the client's view is unchanged,
  which is the entire point, and the teams no longer share a deploy.}
  \label{fig:ch01-supergraph-shape}
\end{figure}

Two months after the launch, Apollo moved composition out of the gateway's
startup and into a managed build step against a
registry~\autocite{debergalis2019managed}. This sounds like an operational
detail and is closer to the foundation of everything in Part~VI. Once
composition is a build artifact, an incompatible schema change is something
that fails in a pipeline instead of at three in the morning, which is what
chapter~\ref{ch:ci-cd-for-a-federated-graph} is about.

Federation~2 reached general availability in April
2022~\autocite{prasek2022federation2}. The change that matters conceptually is
that types became genuinely shared rather than owned-and-extended: the
\code{extend} keyword stopped being required, fields could deliberately exist
in more than one subgraph, and the composition engine was rewritten to validate
\enquote{all theoretically possible queries} instead of only checking that the
pieces fit together. That last one is the difference between composition as a
merge and composition as a proof, and chapter~\ref{ch:composition} is where you
will feel it.

\subsection*{What actually happened to stitching}

Federation did not kill stitching, and books that say it did are repeating a
press release. The Guild took over \code{graphql-tools} in the spring of
2020~\autocite{tanrikulu2020tools6} and kept developing it, and their position
is a fair statement of the trade to this day: if you want workflows out of the
box, federation is the tool; otherwise stitching \enquote{keeps you in control
of your full system architecture}~\autocite{macwilliam2021stitching}. Control
over your own gateway code is a real thing to want. It is also work you have to
keep doing.

The current chapter of the story is an attempt to stop having this argument
per vendor. In May 2024 a Composite Schemas subcommittee re-convened under the
GraphQL specification working group, with Apollo, ChilliCream, Google, Graphile,
The Guild, Hasura, and IBM contributing, to standardise the composition model
that each gateway had been reinventing~\autocite{graphql2024composite}. As of
this writing the specification is an actively edited draft with no published
version~\autocite{compositeschemas2026draft}, while Apollo's own specification
has continued to move and sits at v2.15~\autocite{apollo2026versions}.

This book follows Apollo Federation~v2, because that is the specification
HotChocolate's subgraph package implements and the one Cosmo composes and
routes. Chapter~\ref{ch:fusion} builds the same graph on Composite Schemas with
Fusion instead, and chapter~\ref{ch:the-gateway-market-and-the-future} looks at
where the standard appears to be going.
